\chapter[DG foundations and actual clock lifts]{DG foundations, clock lifts and endpoint-sensitive extensions}
\label{dg:chapter}
A derived construction begins with a differential and its cohomology, not
with a proposed physical clock. This chapter develops the DG material of the
original 2053 manuscript, lines1308--1478. Throughout this chapter $k$ is a
field. We use the actual infinite nodal
and theta dessin algebras. The source is preserved unchanged. Its unit-boundary
example and self-extension count require corrections; its higher-product
and physical interpretations require explicit additional data.

\section{Degrees, maps and a decisive contracting homotopy}
We use cohomological grading. A unital DG algebra is a graded algebra
$B=\bigoplus_{n\in\mathbb Z}B^n$ with $d:B^n\to B^{n+1}$,
$d^2=0$ and $d(ab)=d(a)b+(-1)^{|a|}a\,d(b)$ on homogeneous elements.
A DG map preserves multiplication, the unit, degrees and the differential.
It is a quasi-isomorphism when it induces isomorphisms on every cohomology
group. A left DG module has the corresponding Leibniz rule for its action.
The exterior generator's degree must be declared; the notation
$\Lambda^\bullet(V)$ alone does not specify negative degrees.

\begin{proposition}[A unit boundary is contractible]\label{dg:unit}
If $\varepsilon\in B^{-1}$ satisfies $d\varepsilon=1$, the underlying
$k$-complex of every unital left DG $B$-module is contractible. In particular $H^*(B)=0$; it is not a DG
resolution of a nonzero field with zero differential.
\end{proposition}
\begin{proof}
Define $h:M^n\to M^{n-1}$ by $h(m)=\varepsilon m$. Then
$d_Mh(m)=m-\varepsilon d_Mm$, so $d_Mh+hd_M=1$. Every cocycle $m$
is therefore $d_Mh(m)$. For a right module the correct formula is
$h(m)=(-1)^n m\varepsilon$: its two compositions have opposite
$d_Mm\varepsilon$ terms and again sum to $m$. Taking the regular module
proves the algebra assertion. The displayed homotopy is $k$-linear; graded
$B$-linearity would require an additional compatibility, such as graded
centrality of $\varepsilon$. A unital chain map to a nonzero field would
send $d\varepsilon$ both to zero and to one, which is impossible.
\end{proof}
The qualifier is essential for a general noncommutative DG algebra.
Over our actual nodal algebra $A$, let $C=(A\xrightarrow{1}A)$ in degrees
$-1,0$ and $B=\operatorname{End}_A(C)$ with commutator differential.
The inverse-arrow homotopy $\varepsilon$ satisfies $d_B\varepsilon=1$.
But a degree-$-1$ graded $B$-linear endomorphism of $C$ must commute with
the even projection onto $C^0$; its only possible component $C^0\to C^{-1}$
is then zero. Thus $C$ is acyclic and its underlying complex contracts,
without being contractible by a $B$-linear homotopy. Derived and homotopy
categories must not be conflated.

For the source's two-piece example $k\oplus k\varepsilon$, the differential
is the identity $k\varepsilon\to k$. Both cohomology groups vanish.
The auxiliary identity $dh+hd=f$ for a Koszul relation $f$ is useful
instead of assuming nonzero cohomology \cite{StacksKoszulDG}.

\section{The genuine infinite node and its DG replacement}
The clean interval dessin on the sphere gives
$A=k[\xi,\eta]/(\xi\eta)$, with basis
$1,\xi,\xi^2,\ldots,\eta,\eta^2,\ldots$. It remains infinite dimensional.
The variables $\xi,\eta$ here are actual dessin loops; they are distinct
from the auxiliary persistence parameters $x,y$ used later in the book.
Set $R=k[\xi,\eta]$ and consider
\[
 K^0=R,\qquad K^{-1}=R\varepsilon,\qquad K^n=0\quad(n\ne0,-1),
 \qquad d(b\varepsilon)=\xi\eta b.
\]
The product is ordinary coefficient multiplication, with
$\varepsilon^2=0$ and $\varepsilon$ commuting with degree-zero coefficients.
The odd square relation is imposed also in characteristic two.
Proposition~\ref{thm:nodal-dg} proves all Leibniz cases and the cohomology.
Explicitly, writing $p=a+b\varepsilon$ and $q=c+e\varepsilon$ gives
\[
 d(pq)=d(p)q+\varsigma(p)d(q),\quad
 \varsigma(a+b\varepsilon)=a-b\varepsilon,\quad
 h(a+b\varepsilon)=a\varepsilon,\quad dh+hd=(\xi\eta)\,1_K.
\]
This sign formula extends the homogeneous rule linearly; $d$ is not an
ordinary ungraded derivation of this commutative square-zero algebra.

\begin{center}
\begin{tikzcd}[column sep=large]
0\arrow[r]&R\varepsilon\arrow[r,"\xi\eta"]\arrow[d,"0"]&R\arrow[r]\arrow[d,"\pi"]&0\\
0\arrow[r]&0\arrow[r]&A\arrow[r]&0
\end{tikzcd}
\end{center}
The diagram compares cochain degrees $-1,0$. The map $\pi$ is quotient
augmentation, not an inclusion of $A$ into $R$. Its kernel is $(\xi\eta)$,
and multiplication by $\xi\eta$ is injective since $R$ is a domain.
Thus the cohomology comparison is an isomorphism in every degree.
The two ambient terms are free over $R$. They are not asserted to be free
$A$-modules: the nonzero element $\xi\eta$ acts on them, so their action
does not even factor through $A$. This DG replacement does not shorten the
infinite projective resolution of a nodal residue simple over $A$.

For an explicit calculation,
$\xi^3\eta^2=d(\xi^2\eta\varepsilon)$, whereas $\xi^3$ and $\eta^2$
survive in $H^0$. A polynomial has a unique normal form
$a(\xi)+b(\eta)-c$, where $a(0)=b(0)=c$; remove every monomial divisible
by $\xi\eta$. Their linear independence follows by setting each variable
to zero. This identifies the quotient with the polynomial branch pullback.

\section{Completion and the diagonal coefficient ring}
Write $S=k[[\xi,\eta]]$. The same two-term construction over $S$ maps to
\[
 \Lambda=S/(\xi\eta)\cong
 \{(a(\xi),b(\eta)):a(0)=b(0)\}.
\]
Here is a coefficientwise proof, avoiding an unjustified interchange of
completion and arbitrary cohomology. The branch-evaluation map sends
$F$ to $(F(\xi,0),F(0,\eta))$. Every equal-constant pair is the image of
$a(\xi)+b(\eta)-a(0)$. Its kernel consists of series whose coefficients
on both coordinate axes are zero. Such a series is uniquely $\xi\eta G$
by shifting both exponents down one. The same shift proves multiplication
by $\xi\eta$ injective. Consequently the completed Koszul model has
$H^{-1}=0$, $H^0=\Lambda$ and no other cohomology. All these maps are
continuous for the maximal-ideal topologies: evaluation preserves powers,
and the product relation raises order by two. This is the actual completed
surface order already used in the node chapter.

An isomorphism $H^0(B)\cong\Lambda$ by itself is not this result.
For example $B=\Lambda\oplus\Lambda\delta$, $|\delta|=-1$, $d=0$,
$\delta^2=0$, has the same $H^0$ and nonzero $H^{-1}=\Lambda$.
Projection to $\Lambda$ is not a quasi-isomorphism. A DG enhancement must
specify the DG map or a justified zigzag and prove the required cohomology
comparisons. Equivalences of derived DG-module categories use an external
change-of-algebra theorem with its hypotheses; they do not follow from a
single ordinary algebra isomorphism on $H^0$.

\section{A clock that really lifts, and its semilinearity}
For $q\in k^\times$, define $\sigma_q(\xi)=q\xi$,
$\sigma_q(\eta)=q^{-1}\eta$. Substitution is a polynomial algebra
automorphism, inverse to $\sigma_{q^{-1}}$, and fixes $\xi\eta$.
It acts coefficientwise on $K$, fixing $\varepsilon$, so
\[
 \sigma_q d(b\varepsilon)=\sigma_q(\xi\eta b)
 =\xi\eta\sigma_q(b)=d\sigma_q(b\varepsilon).
\]
Multiplicativity follows from the two-piece product formula. The completed
substitution preserves each maximal-ideal power, so extends continuously,
with continuous inverse, to the formal DG model and branch pullback.
It is a semilinear chain automorphism over the ambient coefficient ring:
$\sigma_q(rp)=\sigma_q(r)\sigma_q(p)$, not generally an $R$-linear map.

The diagonal DVR $k[[u]]$, $u=\xi+\eta$, has branches $(F(\xi),F(\eta))$.
The clock sends $u$ to the pair $(q\xi,q^{-1}\eta)$. If this is a diagonal
series $G(u)$, its two linear coefficients must agree; hence $q=q^{-1}$.
Conversely, if $q^2=1$, substitution sends $F(u)$ to $F(qu)$ and preserves
the diagonal subring. It fixes that subring pointwise exactly when $q=1$,
as the image of $u$ shows. In characteristic two these alternatives may
coincide. Thus preserving the relation and order topology does not by itself
preserve a chosen fixed coefficient action.

For $k=\mathbb C$ and $q=e^{i\theta}$, the coefficient of
$\xi^i\eta^j$ acquires phase $q^{i-j}$. A selected finite monomial space
with its declared orthonormal basis then has a unitary diagonal action.
Choosing that finite norm does not construct a canonical Hilbert completion
of the order, a self-adjoint arithmetic generator or an $L$-function
comparison. The positive-transfer and unitary-logarithm distinctions in
the physics chapters still apply. Gluing DG clocks further requires actual
compatible maps and coherences; no gluing theorem is inferred here.

\section{Theta self extensions require endpoints}
The source calls its double-three-cycle example spherical. Keep that
rational theta distinct from the elliptic theta. Both have
$\sigma=(012)(345)$, but their dart involutions are respectively
$\alpha=(03)(15)(24)$ and $(03)(14)(25)$.
Write vertices as the sorted $\alpha$-pairs $0,1,2$; for the elliptic
case these are $\{0,3\},\{1,4\},\{2,5\}$.
The simples $S_i$ are the residue simples of the actual infinite dessin
algebra. The exact forbidden-word resolution of Theorem~\ref{thm:resolution}
has two positive-degree arms. After Hom to a residue simple, every
arrow differential vanishes, so Ext dimension counts the terminal vertices
of these arms, not just arrows leaving the initial vertex.

\begin{proposition}[Two different Ext tables]\label{dg:theta-ext}
For $n=0$ the dimension matrix is the identity. For $n>0$, using right-module
and path-travel conventions of the resolution chapter, the elliptic theta has
\[
 \dim_k\operatorname{Ext}^n(S_i,S_j)
    =2\,\mathbf1_{j=i-n\bmod3}.
\]
For the spherical theta, positive even $n$ give $2$ on the diagonal and zero
elsewhere; odd $n$ give zero on the diagonal and $1$ at each other vertex.
\end{proposition}
\begin{proof}
For the elliptic involution, both initial arrows lower the vertex label by
one modulo three; every next forbidden arrow does the same. Hence both
length-$n$ indexing words end at $i-n$. For the spherical involution the
forbidden dart permutation has cycles $(04),(13),(25)$. Each two-step arm
returns to its initial vertex, while its first step reaches one of the two
other vertices, one arm each. This proves all positive degrees; degree zero
is Hom between pairwise distinct residue simples. Minimality of the exact
resolution supplies the Hom differential vanishing already used above.
\end{proof}
Thus neither example has two-dimensional self $\operatorname{Ext}^1$.
The two self-extension tables are
\[
\begin{array}{c|rrrrrrr}
n&0&1&2&3&4&5&6\\\hline
\text{elliptic}&1&0&0&2&0&0&2\\
\text{spherical}&1&0&2&0&2&0&2
\end{array}
\]
These are module Ext groups, not stable AR arrows, finite-cutoff quotient
Ext groups or an asserted surface-intersection classification. The completed
residue-simple calculation follows the earlier formal-series resolution;
other continuous-module and derived equivalences need their own arguments.

\section{Higher products and a usable derived interface}
For an uncurved cohomological $A_\infty$ algebra, $m_n$ has degree $2-n$.
With the graded tensor sign convention its identities are
\[
 \sum_{r+s+t=n}(-1)^{r+st}
 m_{r+1+t}(1^{\otimes r}\otimes m_s\otimes1^{\otimes t})=0,
 \qquad s\ge1.
\]
When $m_1=0$, the $n=3$ identity makes $m_2$ associative. The operation
$m_3$ is therefore not its ordinary associator. Minimal-model existence
over a field and the appropriate derived comparison are external results
of Keller \cite[\S\S3.1,3.3,6.2]{KellerAinf}; they are not formalized here.
For an ordinary algebra concentrated in degree zero, taking $m_2$ to be
its product and all other operations zero already gives a minimal model.
For our nodal DG resolution the cohomology is likewise concentrated in zero,
so every higher $m_n$, $n\ge3$, has zero target by degree. This model
therefore supplies no compulsory nonzero higher operation or physical
curvature.

The Ext algebra of a sum of simples is a different graded object. To build
its derived composition data, choose an appropriate projective resolution
$P$ and use the DG endomorphisms with
$\partial f=d_Pf-(-1)^{|f|}fd_P$. Expanding twice cancels the mixed terms
and leaves $d_P^2f-fd_P^2=0$. Bounded-above projective resolutions in the
ordinary module category give the familiar derived-Hom model; completed
and topological versions must specify their Hom spaces and admissible maps.

There is also a useful grading test. Suppose an Ext minimal model is diagonal
in cohomological/internal degree and its operations preserve internal weight.
On inputs of cohomological degrees $a_1,\ldots,a_n$, $m_n$ has output
degree $\sum a_i+2-n$ but internal weight $\sum a_i$. Diagonality forces
$n=2$ unless the output vanishes. The polynomial quadratic-monomial theta
resolutions are linear, so their Ext has this diagonal grading. A transfer
using contractions chosen separately in each internal weight over a field
retains it. Thus an internally compatible model can have every higher
$m_n$ zero. This statement concerns that graded polynomial interface; it
is not an identification of all ungraded or completed enhancements, nor
an assertion that an arbitrary numerically fitted operation is a minimal
model. No physical curvature or arithmetic positivity follows from the
symbol $m_3$.


\section{Derived traces and genuine filtered DG modules}
The graded enveloping algebra is $B^e=B\otimes_k B^{\mathrm{op}}$ with
Koszul signs; Hochschild objects are the derived tensor
$B\otimes^{\mathbf L}_{B^e}B$ and derived Hom
$\mathbf R\operatorname{Hom}_{B^e}(B,B)$, with their homological or
cohomological conventions declared. Derived base tensor or flatness is
needed beyond the field setting. Completed continuous Hochschild theory
also requires its completed tensor and continuity conventions. These
objects are not interchangeable with an arbitrary ordinary matrix trace.

For the ordinary commutative nodal algebra, one concrete calculation is
$HH_0(A)=A/[A,A]=A$: the degree-one Hochschild boundary sends
$a\otimes b$ to $ab-ba=0$. This quotient of the infinite-dimensional algebra carries universal trace
classes but chooses neither a functional nor a positive state. Over $\mathbb C$ declare
$\xi,\eta$ self-adjoint and conjugate coefficients under $*$. The
origin evaluation $\varphi(a)=a(0,0)$ is a clock-invariant positive state,
since $\varphi(a^*a)=|a(0,0)|^2$. Its negative is still an invariant trace
but is not positive, as its value at one shows. Even the positive choice
has GNS null space $(\xi,\eta)$, so its GNS Hilbert space is just
$\mathbb C$. This follows directly from the displayed seminorm, and
shows that trace classes, positivity and a useful spectral representation
are three separate choices. No arithmetic determinant is obtained here.

A strict filtered DG module is a functor from a poset: its maps preserve
degrees, actions and differentials, identities are identities, and every
comparable triple has exactly the same composite map. A homotopy-coherent
replacement must supply the higher coherences too; pairwise homotopies
alone do not define it. For an actual example over our nodal DG model,
\[
 F_{a,b}=\xi^a\eta^b K,\qquad (a,b)\in(\mathbb N^2)^{\mathrm{op}},
\]
forms a family of DG $K$-submodules, with inclusion maps when
$a\ge c$ and $b\ge d$. Indeed multiplying by the central degree-zero
monomial commutes with $d$, and these ideals are closed under the $K$-action.
Inclusions compose literally; the clock preserves each ideal up to its
scalar $q^{a-b}$. Componentwise divisibility also gives
$F_{a,b}\cap F_{c,d}=F_{\max(a,c),\max(b,d)}$. This is a decreasing
module filtration indexed with the declared reversed order.

The declared involution also matters. For self-adjoint $\xi,\eta$, a
nonreal phase $q$ makes $\sigma_q$ an algebra automorphism which is not a
$*$-automorphism. If a unitary implements it in a $*$-representation and
$X=\pi(\xi)$ is self-adjoint, then $UXU^*=qX$ is self-adjoint, so
$(q-\overline q)X=0$ and hence $X=0$; the same holds for $\pi(\eta)$.
Thus this particular nonreal clock cannot act faithfully by unitary
conjugation with those self-adjoint branch conventions. This does not
exclude other involutions or graph-operator constructions with different
adjoint generators; it forbids silently identifying their adjoints with
these declared ones.

The finite monomial windows below instead form increasing subcomplexes of
$k$-vector spaces. They are not $R$-submodules or DG $K$-submodules, because
multiplying by a variable leaves the box. Inverting both variables to
manufacture increasing fractional modules would make $f=\xi\eta$ a unit:
then $d(\varepsilon/f)=1$ and the localized DG algebra contracts. Such
localization loses the node rather than improving its derived resolution.

The remaining source interfaces need their own constructions: bar/cobar,
Ginzburg and Fukaya enhancements are not selected by merely matching $H^0$;
gluing up to homotopy requires coherent maps. Nor does a symbol $m_3$
identify a Hochschild deformation or a measured curvature. The existing
unitary-logarithm discussion is retained in the physics chapter. Its
nonunitary counterpart needs care: a singular finite matrix cannot be an
exponential, since every matrix exponential is invertible. An invertible
unipotent matrix has a finite nilpotent logarithm in characteristic zero,
which is not a self-adjoint generator unless the matrix is the identity.
These observations delimit the operator interface without assuming a
physical realization of the algebraic clock.

\section{Executable windows and exact verification coverage}
The canonical laboratory shows a finite subcomplex, not a finite parent
algebra. For $d\varepsilon=\xi\eta$ it retains source monomials
$0\le i,j\le N-1$ and target monomials $0\le i,j\le N$.
Every source maps injectively to the distinct target $(i+1,j+1)$.
The finite rank is $N^2$, negative cohomology vanishes, and the surviving
$2N+1$ target monomials lie on the two axes. These windows exhaust the
polynomial resolution. For $d\varepsilon=1$ both layers instead use the
same box and contract; for $d=0$ both layers survive. The window is not an $R$-submodule and is not
closed under multiplication, so it must not be advertised as a finite DG
subalgebra. The phase weight $i-j$ is equal at the two ends of every arrow.
The second view displays the exact endpoint-sensitive theta Ext matrices.

\texttt{NodalDG.lean} uses the actual square-zero exterior-line coefficient
algebra, with even/odd embeddings, differential and parity sign. It checks
the Leibniz, homotopy, augmentation and semilinear-clock identities, as well
as the exact ambient polynomial sequence for the nodal relation. These are
actual algebra/module components. A categorical $\mathbb Z$-graded DG object,
completed topology, A-infinity transfer and derived equivalence are not
thereby kernel checked. Exact Python diagnostics check symbolic identities,
normal forms and the two actual dart systems. Browser figures provide finite
windows and phase illustrations; they are not numerical proofs of the
infinite statements or of GRH.
